% FIGURE NEEDED: two concentric circular orbits around S with Earth at E1
% (opposition, Sun--E1--M1 aligned) and at E2 (quadrature, right angle at
% E2), Mars at M1 and M2; 2009 theory solutions booklet, p. 53, Solution 13.
We know the orbital period of Mars then the angle $\angle M_1 S M_2$ can be
determined simply
\[
\angle M_1 S M_2 = \frac{106}{687}\times 360 = 55.5^\circ
\]
\hfill(2 points)

By the same way we determine $\angle E_1 S E_2$:
\[
\angle E_1 S E_2 = \frac{106}{365}\times 360 = 104.5^\circ
\]
\hfill(2 points)

Then angle $\angle M_2 S E_2$ is:
\[
\angle M_2 S E_2 = 104.5 - 55.5 = 49.0^\circ
\]
\hfill(2 points)
\[
\cos(\angle M_2 S E_2) = \frac{SE_2}{SM_2}
\]
\hfill(2 points)
\[
r_{\mathrm{mars}} = \frac{SM_2}{SE_2}
 = \frac{1}{\cos(\angle M_2 S E_2)}
 = \frac{1}{\cos(49^\circ)} = 1.52\ \mathrm{AU}
\]
\hfill(2 points)
